% Options for packages loaded elsewhere
\PassOptionsToPackage{unicode}{hyperref}
\PassOptionsToPackage{hyphens}{url}
\PassOptionsToPackage{dvipsnames,svgnames,x11names}{xcolor}
\documentclass[
  11pt,
  a4paper,
]{article}
\usepackage{xcolor}
\usepackage[margin=1in]{geometry}
\usepackage{amsmath,amssymb}
\setcounter{secnumdepth}{-\maxdimen} % remove section numbering
\usepackage{iftex}
\ifPDFTeX
  \usepackage[T1]{fontenc}
  \usepackage[utf8]{inputenc}
  \usepackage{textcomp} % provide euro and other symbols
\else % if luatex or xetex
  \usepackage{unicode-math} % this also loads fontspec
  \defaultfontfeatures{Scale=MatchLowercase}
  \defaultfontfeatures[\rmfamily]{Ligatures=TeX,Scale=1}
\fi
\usepackage{lmodern}
\ifPDFTeX\else
  % xetex/luatex font selection
  \ifLuaTeX
    \usepackage{luaotfload}
    \directlua{luaotfload.add_fallback("mainfontfallback",{
      "NotoSerifBengali:mode=harf"
    })}
  \fi
  \setmainfont[,RawFeature={fallback=mainfontfallback}]{Latin Modern Roman}
  \ifLuaTeX
    \usepackage{luaotfload}
    \directlua{luaotfload.add_fallback("monofontfallback",{
      "NotoSansBengali:mode=harf"
    })}
  \fi
  \setmonofont[,RawFeature={fallback=monofontfallback}]{DejaVu Sans Mono}
\fi
% Use upquote if available, for straight quotes in verbatim environments
\IfFileExists{upquote.sty}{\usepackage{upquote}}{}
\IfFileExists{microtype.sty}{% use microtype if available
  \usepackage[]{microtype}
  \UseMicrotypeSet[protrusion]{basicmath} % disable protrusion for tt fonts
}{}
\makeatletter
\@ifundefined{KOMAClassName}{% if non-KOMA class
  \IfFileExists{parskip.sty}{%
    \usepackage{parskip}
  }{% else
    \setlength{\parindent}{0pt}
    \setlength{\parskip}{6pt plus 2pt minus 1pt}}
}{% if KOMA class
  \KOMAoptions{parskip=half}}
\makeatother
\usepackage{color}
\usepackage{fancyvrb}
\newcommand{\VerbBar}{|}
\newcommand{\VERB}{\Verb[commandchars=\\\{\}]}
\DefineVerbatimEnvironment{Highlighting}{Verbatim}{commandchars=\\\{\}}
% Add ',fontsize=\small' for more characters per line
\usepackage{framed}
\definecolor{shadecolor}{RGB}{248,248,248}
\newenvironment{Shaded}{\begin{snugshade}}{\end{snugshade}}
\newcommand{\AlertTok}[1]{\textcolor[rgb]{0.94,0.16,0.16}{#1}}
\newcommand{\AnnotationTok}[1]{\textcolor[rgb]{0.56,0.35,0.01}{\textbf{\textit{#1}}}}
\newcommand{\AttributeTok}[1]{\textcolor[rgb]{0.13,0.29,0.53}{#1}}
\newcommand{\BaseNTok}[1]{\textcolor[rgb]{0.00,0.00,0.81}{#1}}
\newcommand{\BuiltInTok}[1]{#1}
\newcommand{\CharTok}[1]{\textcolor[rgb]{0.31,0.60,0.02}{#1}}
\newcommand{\CommentTok}[1]{\textcolor[rgb]{0.56,0.35,0.01}{\textit{#1}}}
\newcommand{\CommentVarTok}[1]{\textcolor[rgb]{0.56,0.35,0.01}{\textbf{\textit{#1}}}}
\newcommand{\ConstantTok}[1]{\textcolor[rgb]{0.56,0.35,0.01}{#1}}
\newcommand{\ControlFlowTok}[1]{\textcolor[rgb]{0.13,0.29,0.53}{\textbf{#1}}}
\newcommand{\DataTypeTok}[1]{\textcolor[rgb]{0.13,0.29,0.53}{#1}}
\newcommand{\DecValTok}[1]{\textcolor[rgb]{0.00,0.00,0.81}{#1}}
\newcommand{\DocumentationTok}[1]{\textcolor[rgb]{0.56,0.35,0.01}{\textbf{\textit{#1}}}}
\newcommand{\ErrorTok}[1]{\textcolor[rgb]{0.64,0.00,0.00}{\textbf{#1}}}
\newcommand{\ExtensionTok}[1]{#1}
\newcommand{\FloatTok}[1]{\textcolor[rgb]{0.00,0.00,0.81}{#1}}
\newcommand{\FunctionTok}[1]{\textcolor[rgb]{0.13,0.29,0.53}{\textbf{#1}}}
\newcommand{\ImportTok}[1]{#1}
\newcommand{\InformationTok}[1]{\textcolor[rgb]{0.56,0.35,0.01}{\textbf{\textit{#1}}}}
\newcommand{\KeywordTok}[1]{\textcolor[rgb]{0.13,0.29,0.53}{\textbf{#1}}}
\newcommand{\NormalTok}[1]{#1}
\newcommand{\OperatorTok}[1]{\textcolor[rgb]{0.81,0.36,0.00}{\textbf{#1}}}
\newcommand{\OtherTok}[1]{\textcolor[rgb]{0.56,0.35,0.01}{#1}}
\newcommand{\PreprocessorTok}[1]{\textcolor[rgb]{0.56,0.35,0.01}{\textit{#1}}}
\newcommand{\RegionMarkerTok}[1]{#1}
\newcommand{\SpecialCharTok}[1]{\textcolor[rgb]{0.81,0.36,0.00}{\textbf{#1}}}
\newcommand{\SpecialStringTok}[1]{\textcolor[rgb]{0.31,0.60,0.02}{#1}}
\newcommand{\StringTok}[1]{\textcolor[rgb]{0.31,0.60,0.02}{#1}}
\newcommand{\VariableTok}[1]{\textcolor[rgb]{0.00,0.00,0.00}{#1}}
\newcommand{\VerbatimStringTok}[1]{\textcolor[rgb]{0.31,0.60,0.02}{#1}}
\newcommand{\WarningTok}[1]{\textcolor[rgb]{0.56,0.35,0.01}{\textbf{\textit{#1}}}}
\usepackage{longtable,booktabs,array}
\usepackage{calc} % for calculating minipage widths
% Correct order of tables after \paragraph or \subparagraph
\usepackage{etoolbox}
\makeatletter
\patchcmd\longtable{\par}{\if@noskipsec\mbox{}\fi\par}{}{}
\makeatother
% Allow footnotes in longtable head/foot
\IfFileExists{footnotehyper.sty}{\usepackage{footnotehyper}}{\usepackage{footnote}}
\makesavenoteenv{longtable}
\setlength{\emergencystretch}{3em} % prevent overfull lines
\providecommand{\tightlist}{%
  \setlength{\itemsep}{0pt}\setlength{\parskip}{0pt}}
% paper_preamble.tex — shared LaTeX preamble for the PCMA/SI paper set.
%
% Compiled with lualatex (NOT pdflatex) because several papers contain
% Bengali script: the CLL paper realises one cognitive state into
% English, Bengali and Banglish, so the Bengali output is EVIDENCE and
% cannot be transliterated away without destroying the point being made.
%
% Font selection is NOT done here. pandoc >= 3 has native
% `mainfontfallback` / `monofontfallback` variables that emit the
% luaotfload fallback registration in the correct position relative to
% unicode-math; doing it by hand in this header lands after unicode-math
% and breaks font loading outright. See tools/build_papers.sh.

\usepackage{amsmath}
\usepackage{amssymb}
\usepackage{mathtools}   % \xleftrightarrow etc. used in cll.md §5
\usepackage{booktabs}

% ---------------------------------------------------------------------
% Unicode symbols appearing in PROSE rather than inside $...$.
%
% The papers were written as markdown and read as markdown, so they use
% real characters — "the ω-boost", "α-child", "ΔR coherence" — instead of
% LaTeX macros. Under unicode-math those codepoints are math-mode-only
% and abort the build ("\symbballowed only in math mode"). Mapping each
% to \ensuremath keeps the markdown readable on GitHub AND compilable,
% which is the whole point of keeping one source.
% ---------------------------------------------------------------------
\usepackage{newunicodechar}

% Greek, lower case
\newunicodechar{α}{\ensuremath{\alpha}}
\newunicodechar{β}{\ensuremath{\beta}}
\newunicodechar{γ}{\ensuremath{\gamma}}
\newunicodechar{δ}{\ensuremath{\delta}}
\newunicodechar{ε}{\ensuremath{\varepsilon}}
\newunicodechar{η}{\ensuremath{\eta}}
\newunicodechar{θ}{\ensuremath{\theta}}
\newunicodechar{κ}{\ensuremath{\kappa}}
\newunicodechar{λ}{\ensuremath{\lambda}}
\newunicodechar{μ}{\ensuremath{\mu}}
\newunicodechar{π}{\ensuremath{\pi}}
\newunicodechar{ρ}{\ensuremath{\rho}}
\newunicodechar{σ}{\ensuremath{\sigma}}
\newunicodechar{τ}{\ensuremath{\tau}}
\newunicodechar{φ}{\ensuremath{\varphi}}
\newunicodechar{χ}{\ensuremath{\chi}}
\newunicodechar{ψ}{\ensuremath{\psi}}
\newunicodechar{ω}{\ensuremath{\omega}}
\newunicodechar{ν}{\ensuremath{\nu}}
\newunicodechar{ξ}{\ensuremath{\xi}}
\newunicodechar{ζ}{\ensuremath{\zeta}}
\newunicodechar{ι}{\ensuremath{\iota}}

% Greek, upper case
\newunicodechar{Δ}{\ensuremath{\Delta}}
\newunicodechar{Φ}{\ensuremath{\Phi}}
\newunicodechar{Σ}{\ensuremath{\Sigma}}
\newunicodechar{Ψ}{\ensuremath{\Psi}}
\newunicodechar{Ω}{\ensuremath{\Omega}}
\newunicodechar{Γ}{\ensuremath{\Gamma}}
\newunicodechar{Λ}{\ensuremath{\Lambda}}
\newunicodechar{Θ}{\ensuremath{\Theta}}
\newunicodechar{Π}{\ensuremath{\Pi}}

% Operators and relations
\newunicodechar{×}{\ensuremath{\times}}
\newunicodechar{·}{\ensuremath{\cdot}}
\newunicodechar{±}{\ensuremath{\pm}}
\newunicodechar{−}{\ensuremath{-}}
\newunicodechar{≈}{\ensuremath{\approx}}
\newunicodechar{≤}{\ensuremath{\leq}}
\newunicodechar{≥}{\ensuremath{\geq}}
\newunicodechar{≠}{\ensuremath{\neq}}
\newunicodechar{∈}{\ensuremath{\in}}
\newunicodechar{∂}{\ensuremath{\partial}}
\newunicodechar{∇}{\ensuremath{\nabla}}
\newunicodechar{√}{\ensuremath{\sqrt{}}}
\newunicodechar{∞}{\ensuremath{\infty}}
\newunicodechar{∝}{\ensuremath{\propto}}
\newunicodechar{∑}{\ensuremath{\sum}}
\newunicodechar{∫}{\ensuremath{\int}}
\newunicodechar{ℝ}{\ensuremath{\mathbb{R}}}
\newunicodechar{ℕ}{\ensuremath{\mathbb{N}}}
\newunicodechar{‖}{\ensuremath{\|}}
\newunicodechar{⊕}{\ensuremath{\oplus}}
\newunicodechar{⊗}{\ensuremath{\otimes}}
\newunicodechar{∀}{\ensuremath{\forall}}
\newunicodechar{∃}{\ensuremath{\exists}}
\newunicodechar{¬}{\ensuremath{\neg}}
\newunicodechar{∧}{\ensuremath{\wedge}}
\newunicodechar{∨}{\ensuremath{\vee}}

% Arrows
\newunicodechar{→}{\ensuremath{\rightarrow}}
\newunicodechar{←}{\ensuremath{\leftarrow}}
\newunicodechar{↔}{\ensuremath{\leftrightarrow}}
\newunicodechar{↑}{\ensuremath{\uparrow}}
\newunicodechar{↓}{\ensuremath{\downarrow}}
\newunicodechar{⇒}{\ensuremath{\Rightarrow}}
\newunicodechar{⇐}{\ensuremath{\Leftarrow}}
\newunicodechar{⇔}{\ensuremath{\Leftrightarrow}}

% Super/subscripts written as literal characters
\newunicodechar{²}{\ensuremath{^{2}}}
\newunicodechar{³}{\ensuremath{^{3}}}
\newunicodechar{⁴}{\ensuremath{^{4}}}
\newunicodechar{⁵}{\ensuremath{^{5}}}
\newunicodechar{⁶}{\ensuremath{^{6}}}
\newunicodechar{₀}{\ensuremath{_{0}}}
\newunicodechar{₁}{\ensuremath{_{1}}}
\newunicodechar{₂}{\ensuremath{_{2}}}
\newunicodechar{₃}{\ensuremath{_{3}}}
\newunicodechar{ᵀ}{\ensuremath{^{\mathsf{T}}}}

% Status glyphs. The papers use emoji as status markers in tables
% (implemented / failed / warning). Mapped to typographic equivalents:
% a colour emoji font in lualatex is unreliable, and the MEANING is what
% carries — a reader needs to see "this one failed", not a red cross.
\newunicodechar{✅}{\ensuremath{\checkmark}}
\newunicodechar{✓}{\ensuremath{\checkmark}}
\newunicodechar{❌}{\ensuremath{\times}}
\newunicodechar{✗}{\ensuremath{\times}}
\newunicodechar{⚠}{\textbf{!}}
\newunicodechar{🧩}{\textbf{[piece]}}
\newunicodechar{📋}{\textbf{[list]}}
\newunicodechar{🟡}{\textbf{[partial]}}
\newunicodechar{🔴}{\textbf{[dead]}}
\newunicodechar{🟢}{\textbf{[live]}}
\newunicodechar{▼}{\ensuremath{\blacktriangledown}}
\newunicodechar{►}{\ensuremath{\blacktriangleright}}
\newunicodechar{□}{\ensuremath{\square}}
\newunicodechar{█}{\rule{0.6em}{0.9ex}}
\newunicodechar{•}{\textbullet}
\newunicodechar{†}{\textdagger}
\newunicodechar{‑}{-}
\newunicodechar{‐}{-}
\newunicodechar{…}{\ldots}
\newunicodechar{।}{\textbar}   % danda: rendered by the Bengali fallback in text
\newunicodechar{⁻}{\ensuremath{^{-}}}

% Zero-width formatting codepoints. These are invisible in an editor but
% are real characters to LaTeX: an emoji written as "⚠️" is TWO
% codepoints (U+26A0 U+FE0F), so declaring only the visible one leaves a
% variation selector with no definition and the build dies on a symbol
% nobody can see. Both are mapped to nothing.
\newunicodechar{^^^^fe0f}{}    % VARIATION SELECTOR-16
\newunicodechar{^^^^200d}{}    % ZERO WIDTH JOINER
\newunicodechar{^^^^200b}{}    % ZERO WIDTH SPACE

\usepackage{longtable}
\usepackage{array}

% Comparison matrices and cost ledgers are wide. Let them shrink to the
% text block rather than overflow the margin.
\usepackage{adjustbox}

% Code blocks: wrap long lines instead of running off the page. Several
% papers quote real C++ with long identifiers.
\usepackage{fvextra}
\fvset{breaklines=true,breakanywhere=true,fontsize=\small}

\setlength{\parskip}{0.5em}
\setlength{\parindent}{0pt}

\usepackage{fancyhdr}
\pagestyle{fancy}
\fancyhf{}
\fancyfoot[C]{\thepage}
\fancyhead[L]{\small\itshape Synthetic Intelligence — research lineage}
\fancyhead[R]{\small\itshape irfan@furylogic.com}
\renewcommand{\headrulewidth}{0.2pt}
\setlength{\headheight}{14pt}
\setlength{\headsep}{18pt}
\usepackage{bookmark}
\IfFileExists{xurl.sty}{\usepackage{xurl}}{} % add URL line breaks if available
\urlstyle{same}
\hypersetup{
  pdftitle={Reasoning and Discovery --- Derivation, Abduction and Runtime Acquisition},
  pdfauthor={Irfan Mahir --- Independent Researcher, Dhaka, Bangladesh --- irfan@furylogic.com},
  colorlinks=true,
  linkcolor={black},
  filecolor={Maroon},
  citecolor={Blue},
  urlcolor={Blue},
  pdfcreator={LaTeX via pandoc}}

\title{Reasoning and Discovery --- Derivation, Abduction and Runtime
Acquisition}
\author{Irfan Mahir --- Independent Researcher, Dhaka, Bangladesh ---
irfan@furylogic.com}
\date{}

\begin{document}
\maketitle

{
\hypersetup{linkcolor=}
\setcounter{tocdepth}{2}
\tableofcontents
}
\section{Reasoning and Discovery --- Derivation, Abduction and Runtime
Acquisition}\label{reasoning-and-discovery-derivation-abduction-and-runtime-acquisition}

\textbf{Irfan Mahir}\\
\emph{Independent Researcher, Dhaka, Bangladesh}\\
irfan@furylogic.com

\begin{quote}
Proof-passes over a truth graph, handling the unknown, and acquiring
knowledge while running.
\end{quote}

\begin{center}\rule{0.5\linewidth}{0.5pt}\end{center}

\subsection{Contents of This Volume}\label{contents-of-this-volume}

\begin{longtable}[]{@{}
  >{\raggedright\arraybackslash}p{(\linewidth - 4\tabcolsep) * \real{0.3333}}
  >{\raggedright\arraybackslash}p{(\linewidth - 4\tabcolsep) * \real{0.3333}}
  >{\raggedright\arraybackslash}p{(\linewidth - 4\tabcolsep) * \real{0.3333}}@{}}
\toprule\noalign{}
\begin{minipage}[b]{\linewidth}\raggedright
\#
\end{minipage} & \begin{minipage}[b]{\linewidth}\raggedright
Paper
\end{minipage} & \begin{minipage}[b]{\linewidth}\raggedright
File
\end{minipage} \\
\midrule\noalign{}
\endhead
\bottomrule\noalign{}
\endlastfoot
1 & RADE: The Recursive Axiom-Derivation Engine --- Deduction over
Prediction in a Persistent Truth Graph & \texttt{RADE.md} \\
2 & Morphological Abduction Architecture (MAA): Overcoming Topological
Local Minima in Graph-Rewriting Cognition Engines & \texttt{maa.md} \\
3 & The L9-Cognitive Graph: Nine Layers from Lexical Decomposition to
Paradigm Generation & \texttt{l9.md} \\
4 & LTO: Live Training Over Manifolds --- A Non-Transformers Continuous
Learning Architecture & \texttt{lto.md} \\
\end{longtable}

\begin{center}\rule{0.5\linewidth}{0.5pt}\end{center}

\newpage

\subsection{RADE: The Recursive Axiom-Derivation Engine --- Deduction
over Prediction in a Persistent Truth
Graph}\label{rade-the-recursive-axiom-derivation-engine-deduction-over-prediction-in-a-persistent-truth-graph}

\begin{quote}
\emph{``Novelty by guessing and novelty by formal expansion look alike
from the outside. Only one of them can show its work.''}
\end{quote}

\begin{center}\rule{0.5\linewidth}{0.5pt}\end{center}

\subsubsection{Abstract}\label{abstract}

A Transformer-style attention mechanism remains a statistical prediction
machine regardless of how its routing is rearranged: what is computed is
a distribution over continuations conditioned on co-occurrence, and
changing the plumbing does not change that. This paper specifies
\textbf{RADE}, an architecture organised around \textbf{logical
synthesis} rather than probabilistic approximation.

RADE replaces three components at once. Knowledge moves from
floating-point weights into an \textbf{Axiom Substrate} --- a persistent
graph of verified relations. Inference moves from a forward pass to a
\textbf{proof pass} over that graph. And termination moves from a fixed
depth to \textbf{homeostatic inhibition}, a resource budget whose
exhaustion triggers a lateral switch to a different region of the graph
rather than an unbounded search.

The consequence that matters is a change of failure mode. A predictive
system fails by producing fluent falsehood, which is undetectable from
the output alone. RADE fails by \textbf{logical inconsistency}, which is
detectable, localisable, and attributable to specific premises.

\begin{center}\rule{0.5\linewidth}{0.5pt}\end{center}

\subsubsection{1. The Problem with Rearranged
Plumbing}\label{the-problem-with-rearranged-plumbing}

Architectural variation within the predictive family --- sparse
attention, linear attention, state-space substitutions, mixture routing
--- changes the cost of computing \(P(w_n \mid w_1 \ldots w_{n-1})\)
without changing what that quantity is. A count of how often one symbol
followed another is not an account of \emph{why} it did, and no
rearrangement of the machinery that computes the count converts it into
one.

This is not an argument that prediction is useless. It is an argument
about what a system can be asked for. A predictive system can be asked
for a continuation. It cannot be asked for a \emph{derivation}, because
it does not have one to give --- the output was not produced by steps
that could be exhibited.

RADE is specified to answer the second kind of request.

\begin{center}\rule{0.5\linewidth}{0.5pt}\end{center}

\subsubsection{2. The Axiom Substrate}\label{the-axiom-substrate}

\paragraph{2.1 Knowledge as structure, not as
parameters}\label{knowledge-as-structure-not-as-parameters}

In a Transformer, knowledge is distributed across floating-point weights
and cannot be inspected, edited, or deleted at the granularity of a
fact. In RADE it is held explicitly:

\begin{itemize}
\tightlist
\item
  \textbf{Formal objects} \(O\) --- concepts, as graph vertices.
\item
  \textbf{Operators} \(\Phi\) --- typed relations, as directed edges.
\end{itemize}

The distinction from embeddings is persistence semantics. An embedding
degrades under continued training: the representation of a concept
drifts as unrelated updates accumulate. A graph edge does not. If the
substrate holds \texttt{A\ is\_a\ B}, it holds it until something
explicitly removes it, and the removal is an event that can be logged.

\paragraph{2.2 Why this matters for
correction}\label{why-this-matters-for-correction}

A predictive model cannot be told that one specific thing it believes is
wrong. The belief is not localised, so correction is retraining, and
retraining perturbs everything else.

In a truth graph, a wrong fact is an edge. Removing it is \(O(1)\), and
the consequences of removing it are computable: every derived edge
citing it as a premise becomes unsupported and can be retracted by
following the provenance chain. Correction is therefore \emph{bounded}
rather than global.

\begin{center}\rule{0.5\linewidth}{0.5pt}\end{center}

\subsubsection{3. The Proof Pass}\label{the-proof-pass}

The forward pass is replaced by a traversal.

\textbf{Input.} A query --- for example, \emph{why is X a
superconductor?}

\textbf{Operation.} The engine does not estimate an answer. It searches
the graph for a path connecting the queried concepts through typed
operators. If a path exists, the path \emph{is} the answer, and every
edge along it is either an asserted premise or a derived edge that names
the two premises that produced it.

\textbf{When no path exists.} The engine enters a
\textbf{search-and-constrain} phase: it proposes a candidate link and
tests it against the existing axioms. A candidate that contradicts a
held axiom is rejected. A candidate that survives is committed as a
derived edge, with its provenance recorded.

This is the distinction the paper turns on:

\begin{quote}
\textbf{Novelty by guessing} produces a plausible string.
\textbf{Novelty by formal expansion} produces an edge that survived a
consistency test and can name what it rests on.
\end{quote}

Both look like new output. Only the second can be audited, and only the
second can be retracted cleanly when its support is withdrawn.

\begin{center}\rule{0.5\linewidth}{0.5pt}\end{center}

\subsubsection{4. Homeostatic Inhibition}\label{homeostatic-inhibition}

An unbounded proof search does not terminate on a large graph. The
biological analogue of the fix is inhibitory interneurons, which prevent
excitation from propagating without limit.

\textbf{Resource budget.} Every search path is allocated a bounded
budget.

\textbf{Lateral inhibition.} A path that yields no result within its
budget is abandoned, and the search is displaced to a different region
of the graph rather than deepened.

The behavioural effect is the ability to \emph{stop pursuing an
unproductive line}, which is a capability rather than a limitation: a
system that cannot abandon a path cannot allocate effort. This mechanism
is the direct ancestor of the tension-bounded construction budget in the
successor architecture, where the budget is derived from a physical
observable rather than set as a constant.

\begin{center}\rule{0.5\linewidth}{0.5pt}\end{center}

\subsubsection{5. Comparison}\label{comparison}

\begin{longtable}[]{@{}
  >{\raggedright\arraybackslash}p{(\linewidth - 4\tabcolsep) * \real{0.3333}}
  >{\raggedright\arraybackslash}p{(\linewidth - 4\tabcolsep) * \real{0.3333}}
  >{\raggedright\arraybackslash}p{(\linewidth - 4\tabcolsep) * \real{0.3333}}@{}}
\toprule\noalign{}
\begin{minipage}[b]{\linewidth}\raggedright
Property
\end{minipage} & \begin{minipage}[b]{\linewidth}\raggedright
Transformer
\end{minipage} & \begin{minipage}[b]{\linewidth}\raggedright
RADE
\end{minipage} \\
\midrule\noalign{}
\endhead
\bottomrule\noalign{}
\endlastfoot
Foundation & statistical probabilities & logical axioms \\
Inference & forward pass & proof pass over a graph \\
Novelty & interpolation, indistinguishable from hallucination & formal
expansion, consistency-tested \\
Memory & lossy KV cache & persistent graph of hard relations \\
Correction & retraining, global & edge removal, bounded by provenance \\
Failure mode & fluent falsehood, \textbf{undetectable from output} &
logical inconsistency, \textbf{detectable and localisable} \\
Objective & minimise prediction error & maximise logical consistency \\
\end{longtable}

The failure-mode row carries the argument. A system whose errors are
detectable can be improved by an operator who does not have access to
its internals. A system whose errors are fluent cannot.

\begin{center}\rule{0.5\linewidth}{0.5pt}\end{center}

\subsubsection{6. Ingestion: The
Axiom-ifier}\label{ingestion-the-axiom-ifier}

A corpus is not knowledge until it is structured. RADE separates
\textbf{reading} from \textbf{reasoning}: an ingestion stage converts
source material into structured triples, and the reasoning engine never
touches raw text.

The ingestion stage emits only triples that the substrate can type-check
against its operator set. Anything it cannot structure is dropped rather
than approximated --- a coverage gap is preferable to a fabricated
relation, because a gap is visible and a fabrication is not.

The successor architecture implements exactly this split: mechanical AST
extraction produces triples from source, and every fact carries a
provenance class distinguishing what was asserted from what was derived
and what was established by experiment.

\begin{center}\rule{0.5\linewidth}{0.5pt}\end{center}

\subsubsection{7. Status and Relation to the Implemented
System}\label{status-and-relation-to-the-implemented-system}

RADE is an \textbf{architectural proposal}, written before the substrate
that now exists. Its components map onto implemented mechanisms as
follows:

\begin{longtable}[]{@{}
  >{\raggedright\arraybackslash}p{(\linewidth - 4\tabcolsep) * \real{0.3333}}
  >{\raggedright\arraybackslash}p{(\linewidth - 4\tabcolsep) * \real{0.3333}}
  >{\raggedright\arraybackslash}p{(\linewidth - 4\tabcolsep) * \real{0.3333}}@{}}
\toprule\noalign{}
\begin{minipage}[b]{\linewidth}\raggedright
RADE component
\end{minipage} & \begin{minipage}[b]{\linewidth}\raggedright
Implemented as
\end{minipage} & \begin{minipage}[b]{\linewidth}\raggedright
Status
\end{minipage} \\
\midrule\noalign{}
\endhead
\bottomrule\noalign{}
\endlastfoot
Axiom Substrate & typed concept graph, \texttt{data/*.axioms} &
\textbf{live} \\
Proof pass & forward chaining with premise citation & \textbf{live} ---
derived edges name their parents \\
Search-and-constrain & contradiction check before assertion &
\textbf{live} \\
Resource budget & tension-bounded construction budget & \textbf{live},
now derived from field energy \\
Lateral inhibition & budget exhaustion ends a derivation chain &
\textbf{partial} --- no region-switching heuristic \\
Axiom-ifier & mechanical AST extraction & \textbf{live}, no ML \\
\end{longtable}

The one component not carried across is the explicit region-switching
behaviour of §4. The current engine exhausts a budget and stops; it does
not redirect to a different subgraph. Whether that redirection is worth
implementing is open, and it should not be reported as present.

\begin{center}\rule{0.5\linewidth}{0.5pt}\end{center}

\subsubsection{8. Limitations}\label{limitations}

\textbf{Ingestion is the bottleneck, not reasoning.} The proof pass is
cheap; converting unstructured source into typed triples is not. The
quality ceiling of the whole system is the quality of that conversion,
and this paper does not solve it.

\textbf{Consistency is not truth.} The search-and-constrain phase
rejects candidates that contradict held axioms. It cannot reject a
candidate that is consistent with a set of axioms that are themselves
wrong. Consistency is a necessary condition and nothing more.

\textbf{Coverage cliff.} A query whose concepts are absent from the
substrate returns nothing. This is the correct behaviour and it is also
a visible limitation; the honest presentation is an explicit gap rather
than a degraded guess.

\begin{center}\rule{0.5\linewidth}{0.5pt}\end{center}

\subsubsection{9. Conclusion}\label{conclusion}

RADE's claim is narrow and, stated carefully, defensible: an
architecture that stores relations rather than parameters, infers by
traversal rather than by prediction, and bounds its search by an
explicit budget can answer for its outputs in a way a predictive system
structurally cannot.

It does not claim to reason better. It claims to reason \emph{auditably}
--- and to fail in a way that can be seen.

\begin{center}\rule{0.5\linewidth}{0.5pt}\end{center}

\subsubsection{References}\label{references}

\begin{itemize}
\tightlist
\item
  Mahir, I. \emph{From Artificial to Synthetic Intelligence: Principles
  of Constructive Structural Cognition.} --- the tri-layer architecture
  in which RADE's components are implemented.
\item
  Mahir, I. \emph{JAS: Einstein's Epistemology as a Running
  Architecture.} --- the provenance classes that make retraction
  bounded.
\item
  Implementation: \texttt{src/rade.hpp}.
\end{itemize}

\newpage

\subsection{Morphological Abduction Architecture (MAA): Overcoming
Topological Local Minima in Graph-Rewriting Cognition
Engines}\label{morphological-abduction-architecture-maa-overcoming-topological-local-minima-in-graph-rewriting-cognition-engines}

\textbf{Author:} Irfan Mahir \textbf{Date:} July 2026 \textbf{Subjects:}
Artificial Intelligence (cs.AI); Computation and Language (cs.CL); Logic
in Computer Science (cs.LO)

\begin{center}\rule{0.5\linewidth}{0.5pt}\end{center}

\subsubsection{Abstract}\label{abstract-1}

Graph-rewriting cognition engines offer a rigid, verifiable alternative
to the fluid approximations of statistical large language models.
However, when these architectures treat semantic inputs as atomic,
unbreakable nodes, they become highly susceptible to topological
stagnation. When encountering an undefined concept (an orphan node),
standard inference pipelines (deduction, induction, analogy) fail
because the search space is mathematically empty, causing the system to
infinitely loop its active cache. We propose the Morphological Abduction
Architecture (MAA), a framework that breaks the atomic constraint by
introducing Sub-Lexical Decomposition and Abductive Synthesis. By
dynamically fracturing unknown nodes and mapping their sub-components to
existing ontologies, the system is able to hypothesize structural
meaning---transforming the engine from a rigid memorizer into a dynamic
semantic synthesizer capable of state-dependent cognitive dispatching.

\begin{center}\rule{0.5\linewidth}{0.5pt}\end{center}

\subsubsection{1. Introduction}\label{introduction}

In discrete graph-rewriting systems, cognition is defined by the
traversal and construction of edges between known concept nodes. While
this ensures strict logical adherence, it introduces a critical
vulnerability: the ``Atomic Node Trap.''

When standard graph engines encounter a novel term---such as a
previously unseen acronym or portmanteau---the parser initializes a
singular, isolated node \(\emptyset_V\). Because this node possesses a
degree of zero (\(deg(\emptyset_V) = 0\)), there are no topological
pathways for the inference engine to traverse. Consequently, monolithic
evaluation loops (which linearly execute deduce \(\rightarrow\) induce
\(\rightarrow\) analogize pipelines) encounter a void. Lacking a
mechanism to process the unknown, the engine falls back onto its
highest-activation localized structures, resulting in pathological
repetition (e.g., infinitely looping a cached analogy about
consciousness regardless of the prompt).

To achieve general reasoning, a system must be capable of synthesizing
meaning from the unknown. The Morphological Abduction Architecture (MAA)
formalizes the mechanics of ``connecting the dots'' by physically
fragmenting orphan nodes and generating hypothesized graph structures.

\subsubsection{2. Component 1: Morphological
Decomposition}\label{component-1-morphological-decomposition}

To process an unknown variable, MAA dictates that a concept node must
not be treated as a fundamental particle, but as a composite structure
capable of fission.

Let the user introduce an unknown string, \(C_{unk}\). When the parser
detects that \(deg(C_{unk}) = 0\) within the active ontology \(O\), the
engine suspends standard inference and triggers a Decomposition Rule.

The string \(C_{unk}\) is fractured into a set of sub-lexical fragments
\(F = \{f_1, f_2, \dots, f_n\}\) based on linguistic boundaries,
syllables, or character clusters. The engine then projects these
fragments against the existing ontology to identify partial topological
matches:

\[ M = \{\, f_i \in F \mid \exists n \in O \text{ where } similarity(f_i, n) > \tau \,\} \]

\textbf{Case Study (The Superuser Maneuver):} If the engine encounters
the unknown node \texttt{{[}sudo{]}}, it evaluates to
\(deg([sudo]) = 0\). MAA fractures the string into \(F = \{su, do\}\).
The engine scans the ontology and successfully binds the fragments to
known concepts: \(su \rightarrow [super\_user]\)
\(do \rightarrow [execute\_action]\)

The engine has successfully extracted topological anchors from a
previously isolated node.

\subsubsection{3. Component 2: Abductive Synthesis
Graphing}\label{component-2-abductive-synthesis-graphing}

Once morphological anchors are established, the engine cannot use
Deduction (which requires absolute axiomatic truth) or Induction (which
requires statistical precedence). It must use \textbf{Abduction}---the
generation of the most likely explanatory hypothesis.

MAA introduces a volatile edge type to the graph schema:
\texttt{EDGE\_HYPOTHESIZED}. Unlike axiomatic edges, hypothesized edges
are temporary structural bridges used to test theoretical topologies
without corrupting the permanent ontology.

Following the extraction of match set \(M\), the Abductive Synthesis
Engine generates a Hypothesis Graph \(\Gamma_H\).

For the unknown node \texttt{{[}rdma{]}}, the engine fragments the node
into \(\{rd, ma\}\). It finds partial matches for
\(rd \rightarrow [remote / read]\) and
\(ma \rightarrow [memory / machine]\). The engine synthesizes a cluster
of theoretical connections:
\texttt{{[}rdma{]}\ -\/-EDGE\_HYPOTHESIZED-\/-\textgreater{}\ {[}remote\_memory{]}}
\texttt{{[}rdma{]}\ -\/-EDGE\_HYPOTHESIZED-\/-\textgreater{}\ {[}read\_machine{]}}

By generating these volatile edges, the node \(C_{unk}\) is no longer
orphaned. The search space is populated, allowing the cognitive loop to
resume execution, score the hypotheses against the surrounding context,
and select the highest-probability path.

\subsubsection{4. State-Dependent Cognitive
Dispatching}\label{state-dependent-cognitive-dispatching}

The final requirement of MAA is the abolition of the monolithic
execution loop. A static \texttt{respond()} pipeline that applies the
same sequence of rules to every input is fundamentally a memorization
behavior. True cognition requires the processing architecture to adapt
to the topology of the current graph state.

MAA replaces the static loop with a dynamic Cognitive Dispatcher
function, \(\Delta\). Before any inference is executed, \(\Delta\)
evaluates the macroscopic geometry of the active graph and routes the
execution thread to specialized sub-engines:

\begin{enumerate}
\def\labelenumi{\arabic{enumi}.}
\tightlist
\item
  \textbf{State \(\alpha\) (Topological Collision):} If the input graph
  introduces an edge that directly opposes an axiomatic invariant (e.g.,
  \texttt{machine\ -\/-has-\/-\textgreater{}\ biology}), \(\Delta\)
  routes to the \textbf{Resolution Pipeline} to isolate and reject the
  contradiction.
\item
  \textbf{State \(\beta\) (Orphan Detected):} If the input graph
  contains a node where \(deg(n) = 0\) (e.g., \texttt{{[}rdma{]}}),
  \(\Delta\) routes to the \textbf{Abductive Pipeline} (Decomposition
  and Synthesis) to map the unknown.
\item
  \textbf{State \(\gamma\) (Connected Query):} Only when the input graph
  is fully integrated into the existing ontology does \(\Delta\) route
  to the \textbf{Standard Inference Pipeline} to deduce new paths.
\end{enumerate}

\[
\Delta(\Gamma_{active}) \rightarrow
\begin{cases}
Pipeline_{Resolution} & \text{if } \exists \text{ contradiction in } \Gamma \\
Pipeline_{Abductive}  & \text{if } \exists n \in \Gamma \text{ where } deg(n) = 0 \\
Pipeline_{Standard}   & \text{otherwise}
\end{cases}
\]

\subsubsection{5. Conclusion}\label{conclusion-1}

The Morphological Abduction Architecture demonstrates that the inability
of discrete graph engines to handle unknown variables is not a
limitation of discrete mathematics, but a failure of parsing granularity
and execution routing. By allowing the engine to fracture concepts into
their morphological roots and abductively synthesize hypothesized edge
relationships, we eliminate topological local minima. Furthermore,
implementing a State-Dependent Cognitive Dispatcher ensures that the
machine actively ``reads'' its own structural state before acting,
permanently transitioning the architecture from a rigid data-retrieval
script into an adaptive, synthesizing cognition engine.

\newpage

\subsection{The L9-Cognitive Graph: Nine Layers from Lexical
Decomposition to Paradigm
Generation}\label{the-l9-cognitive-graph-nine-layers-from-lexical-decomposition-to-paradigm-generation}

\begin{quote}
\emph{``A graph that answers
\texttt{{[}Apple{]}\ -\/-is\_a-\/-\textgreater{}\ {[}Fruit{]}} is a
dictionary. Maxwell did not consult a dictionary.''}
\end{quote}

\begin{center}\rule{0.5\linewidth}{0.5pt}\end{center}

\subsubsection{Abstract}\label{abstract-2}

Probabilistic next-symbol prediction confines a system to interpolating
within existing knowledge. We propose the \textbf{L9-Cognitive Graph}, a
neuro-symbolic architecture organised as nine discrete layers of
traversal logic over a deterministic symbolic graph, ascending from
lexical decomposition to the generation of new axiomatic roots.

The organising claim is that breakthrough reasoning is not semantic
retrieval but \textbf{cross-domain isomorphism}: recognising that the
mathematical structure of one domain describes the unresolved anomalies
of another. Achieving this requires the engine to move from
\emph{constraint enforcement} to \emph{abductive constraint generation}.

This paper specifies all nine layers --- each with its input, operation,
output and failure mode --- rather than listing them. Section 5 states,
layer by layer, which are implemented, which are partial, and which
remain proposals. Four are live, two are partial, and three are not
built.

\begin{center}\rule{0.5\linewidth}{0.5pt}\end{center}

\subsubsection{1. The Limits of Semantic
Retrieval}\label{the-limits-of-semantic-retrieval}

A graph database mapping
\texttt{{[}Apple{]}\ -\/-is\_a-\/-\textgreater{}\ {[}Fruit{]}} is a
semantic dictionary: it returns what was stored, structured. Human
breakthrough logic does not work this way. Maxwell's unification of
electricity and magnetism was not a lookup; it was the recognition that
one domain's mathematical structure resolved another's anomalies.

Reproducing that requires two capabilities a dictionary lacks. The
engine must be able to notice that its own knowledge is
\textbf{structurally inconsistent} rather than merely incomplete, and it
must be able to \textbf{propose structure that was never asserted} in
response.

The layers below are ordered by how far each departs from retrieval.

\begin{center}\rule{0.5\linewidth}{0.5pt}\end{center}

\subsubsection{2. Biological Motivation: Two
Topologies}\label{biological-motivation-two-topologies}

The architecture borrows two structures rather than a metaphor.

\textbf{Neocortical columns --- the local graph.} The brain processes
localised, highly deterministic rules: syntax, basic physics. These
correspond to densely connected subgraph clusters with near-unit edge
weights, where traversal is reliable and cheap.

\textbf{The Default Mode Network --- the abductive engine.} The DMN
activates during undirected thought and deep problem-solving, searching
for low-probability, long-distance connections across disparate columns.
It corresponds to an asynchronous background process running
bridge-finding across structurally similar but semantically distant
subgraphs.

The functional distinction that matters: \textbf{the local graph
answers, the DMN wanders.} A system with only the first can never
surprise itself.

\begin{center}\rule{0.5\linewidth}{0.5pt}\end{center}

\subsubsection{3. The Nine Layers}\label{the-nine-layers}

Each layer is specified by what it takes, what it does, what it emits,
and how it fails. Layers are cumulative: L\(n\) presupposes
L\(1..n{-}1\).

\paragraph{L1 --- Lexical / Morphological
Decomposition}\label{l1-lexical-morphological-decomposition}

\textbf{Input.} A token that resolves to no node. \textbf{Operation.}
Fracture the string into sub-lexical fragments --- morphemes, syllables,
character clusters --- and project each against the ontology for partial
matches. \texttt{biology} decomposes to \texttt{bio} + \texttt{logy};
\texttt{sudo} to \texttt{su} + \texttt{do}. \textbf{Output.} A set of
candidate bindings for an otherwise orphan node. \textbf{Failure mode.}
Spurious decomposition. Short fragments match everything, so a fragment
carries usable evidence only above a minimum length --- an information
floor, not a stopword list.

This layer exists because a node of degree zero is unreachable, and an
unreachable node makes every higher layer inapplicable.

\paragraph{L2 --- Ontological Retrieval}\label{l2-ontological-retrieval}

\textbf{Input.} A resolved node. \textbf{Operation.} Traverse explicit
edges. \texttt{Darwin\ -\textgreater{}\ Human}. \textbf{Output.}
Asserted neighbours. \textbf{Failure mode.} Silence when no edge exists.
This is the honest failure and is preferable to a fabricated one.

L1 and L2 together constitute the dictionary. Everything above is the
argument of this paper.

\paragraph{L3 --- Transitive Inference}\label{l3-transitive-inference}

\textbf{Input.} Two chained edges \(A \to B\), \(B \to C\).
\textbf{Operation.} Instantiate the implied edge \(A \to C\), recording
both premises. \textbf{Output.} A derived edge, distinguishable from an
asserted one. \textbf{Failure mode.} Unbounded expansion. Transitive
closure over a large graph is quadratic, so derivation must be bounded
by an explicit budget.

The provenance requirement is not decoration. A derived edge that cannot
name its premises cannot be retracted when one of them is withdrawn.

\paragraph{L4 --- Causal Abduction}\label{l4-causal-abduction}

\textbf{Input.} An observed end state \(Y\). \textbf{Operation.} Search
\emph{backwards} for a cause \(X\) whose consequences include \(Y\).
This inverts L3's direction of travel. \textbf{Output.} A hypothesis,
not a conclusion. \textbf{Failure mode.} Affirming the consequent. Many
causes produce the same effect, so an abduced \(X\) is a candidate
requiring independent support --- and labelling it otherwise is the
single easiest way for this architecture to produce confident nonsense.

\paragraph{L5 --- Counterfactual
Simulation}\label{l5-counterfactual-simulation}

\textbf{Input.} A core axiom. \textbf{Operation.} Temporarily invert it
--- \emph{what if gravity repelled?} --- and propagate consequences,
checking whether the resulting system remains consistent.
\textbf{Output.} A stability verdict on the inverted world.
\textbf{Failure mode.} Combinatorial explosion, and contamination of the
base graph if the inversion is not perfectly isolated and reverted.

L5 is how the engine determines which axioms are \emph{load-bearing}: an
axiom whose inversion changes nothing was doing no work.

\paragraph{L6 --- Analogical Mapping}\label{l6-analogical-mapping}

\textbf{Input.} Two subgraphs from different domains.
\textbf{Operation.} Attempt a structure-preserving \(1{:}1\)
correspondence --- fluid dynamics onto electrical circuits, pressure
\(\to\) voltage, flow \(\to\) current. \textbf{Output.} A mapping, plus
the edges of one domain transferred onto the other as hypotheses.
\textbf{Failure mode.} Superficial isomorphism. Two subgraphs can share
shape while the correspondence means nothing, so a mapping must be
scored by more than topology.

\textbf{This is the first layer of genuine novelty}, because its output
is not entailed by either domain alone.

\paragraph{L7 --- Anomaly Detection}\label{l7-anomaly-detection}

\textbf{Input.} The graph. \textbf{Operation.} Scan for nodes of high
connectivity carrying \emph{contradictory} constraints --- a concept
that participates widely while holding incompatible commitments. This is
the \textbf{Maxwell gap}: the signature of a theory that is incomplete
at a specific, locatable point. \textbf{Output.} A ranked list of
structural contradictions. \textbf{Failure mode.} Mistaking noise for
anomaly. Extraction error produces contradictions that indicate a bad
parse rather than a gap in the world.

L7 is the layer that lets the system be \emph{dissatisfied} with what it
knows, which is the precondition for L8 and L9.

\paragraph{L8 --- Synthesis (The DMN
Layer)}\label{l8-synthesis-the-dmn-layer}

\textbf{Input.} An L7 anomaly. \textbf{Operation.} Generate a
\textbf{bridge node} resolving it, constrained to satisfy both of the
subgraphs it joins. \textbf{Output.} A new vertex that appeared in no
source. \textbf{Failure mode.} Unfalsifiable patching. A bridge invented
purely to absorb a contradiction, with no independent consequences,
explains nothing --- it renames the problem.

The discipline required here is that a synthesised node must make
predictions beyond the anomaly that motivated it.

\paragraph{L9 --- Paradigm Generation}\label{l9-paradigm-generation}

\textbf{Input.} A region of the graph that is verbose, irregular and
patched. \textbf{Operation.} Define a new \textbf{axiomatic root} from
which the messy region follows, collapsing many special cases into fewer
general ones. \textbf{Output.} A reorganisation of the knowledge base,
not an addition to it. \textbf{Failure mode.} Losing content. A
simplification that fails to entail what it replaced is not a paradigm;
it is deletion.

L9 is the only layer whose success criterion is \emph{compression}: the
same commitments held in less structure.

\begin{center}\rule{0.5\linewidth}{0.5pt}\end{center}

\subsubsection{4. Scoring Novelty}\label{scoring-novelty}

Novelty is not randomness --- it is the discovery of hidden structural
isomorphism. Given subgraphs \(A\) and \(B\), we seek pairs that are
\textbf{semantically distant} but \textbf{structurally identical}:

\[
\Phi(A, B) \;=\; \frac{I(A, B)}{e^{-k \cdot d(A, B)}}
\]

where \(I(A,B)\) is a maximal common subgraph isomorphism score,
\(d(A,B)\) is the shortest path between the subgraph roots, and \(k\) is
a scaling constant.

The exponential in the denominator is the load-bearing choice: it makes
\(\Phi\) grow with distance, so a structural match between
\emph{adjacent} concepts scores poorly. Two neighbouring nodes sharing
shape is expected. Two nodes twelve hops apart sharing shape is a
candidate breakthrough. Pairs exceeding a threshold \(\tau\) are
flagged.

\textbf{A caution this formulation invites.} \(\Phi\) rewards distance
unboundedly, so with a finite isomorphism score the highest-scoring
pairs will tend to be the most remote --- which includes pairs that are
remote because one of them is extraction noise. Any deployment needs a
floor on subgraph quality before distance is rewarded.

\begin{center}\rule{0.5\linewidth}{0.5pt}\end{center}

\subsubsection{5. Implementation Status}\label{implementation-status}

The paper's claims are separated from the engine's behaviour here, layer
by layer.

\begin{longtable}[]{@{}
  >{\raggedright\arraybackslash}p{(\linewidth - 4\tabcolsep) * \real{0.3333}}
  >{\raggedright\arraybackslash}p{(\linewidth - 4\tabcolsep) * \real{0.3333}}
  >{\raggedright\arraybackslash}p{(\linewidth - 4\tabcolsep) * \real{0.3333}}@{}}
\toprule\noalign{}
\begin{minipage}[b]{\linewidth}\raggedright
Layer
\end{minipage} & \begin{minipage}[b]{\linewidth}\raggedright
Status
\end{minipage} & \begin{minipage}[b]{\linewidth}\raggedright
Where, or what is missing
\end{minipage} \\
\midrule\noalign{}
\endhead
\bottomrule\noalign{}
\endlastfoot
\textbf{L1} Morphological & \textbf{live} & bigram sensory lanes; graded
recognition survives misspelling, 7/7 measured \\
\textbf{L2} Retrieval & \textbf{live} & typed concept graph traversal \\
\textbf{L3} Transitive & \textbf{live} & forward chaining with premise
citation; derived edges name both parents \\
\textbf{L4} Abduction & \textbf{partial} & abduction over upstream
cavities exists; no general backward search from an arbitrary \(Y\) \\
\textbf{L5} Counterfactual & \textbf{not built} & requires isolated
graph forking with guaranteed revert; no such isolation exists \\
\textbf{L6} Analogical & \textbf{live} & shared-target intersection plus
rubber-tunnel bridging, with 20-seed reproducibility data; the Fisher
metric supplies a principled distance \\
\textbf{L7} Anomaly & \textbf{partial} & contradiction detection and
isolated-subgraph/logic-gap scans exist; no ranked Maxwell-gap
register \\
\textbf{L8} Synthesis & \textbf{partial} & fission spawns new vertices
to resolve contradiction, and mutation forces distant-pair hypotheses
that decay if unsupported --- but neither is \emph{driven by} an L7
anomaly report \\
\textbf{L9} Paradigm & \textbf{not built} & no mechanism reorganises the
graph around a new root; macro-rule induction compresses
\emph{derivations}, not \emph{axioms}, which is a weaker relative \\
\end{longtable}

\textbf{The honest summary.} The lower half of the stack is real. The
upper half is partially present but \emph{disconnected}: the engine has
anomaly detection and it has synthesis, but synthesis does not consume
anomaly reports. Wiring L7 into L8 is the single highest-value unbuilt
link, and it is engineering rather than research.

L5 and L9 are genuinely open. L9 in particular is not a matter of
effort: no criterion is currently specified for when a proposed
axiomatic root is \emph{better} than the structure it would replace, and
without that criterion the layer cannot be evaluated even if built.

\begin{center}\rule{0.5\linewidth}{0.5pt}\end{center}

\subsubsection{6. Limitations}\label{limitations-1}

\textbf{The layer ordering is asserted.} Nine layers with clean
boundaries is a presentational convenience. In the running engine, L1
recognition and L6 bridging interact directly, and L3 derivation feeds
L7 detection without passing through L4 or L5. The stack is a taxonomy,
not a pipeline.

\textbf{\(\Phi\) is untuned.} The constant \(k\) and threshold \(\tau\)
have no principled setting, and §4 notes the formulation's bias toward
remote noise. The Fisher-information distance is a more defensible
replacement for \(d\), and supersedes the shortest-path form where both
are available.

\textbf{No layer above L3 has an oracle.} L3's output is checkable by
construction. An L6 mapping, an L8 bridge or an L9 paradigm has no
external verdict available --- which is precisely the gap the
compile-verify loop closes for code, and does not close for physics or
biology.

\begin{center}\rule{0.5\linewidth}{0.5pt}\end{center}

\subsubsection{7. Conclusion}\label{conclusion-2}

The L9 stack describes the ascent from a semantic dictionary to an
isomorphic constraint-satisfaction engine --- one able to observe
structural contradictions in what it holds and calculate the structure
required to unify them.

Four layers are live and two are partial. The most useful outcome of
writing the stack out in full is not the top of it but the middle: it
makes visible that anomaly detection and synthesis both exist and are
\textbf{not connected to each other}, which was not apparent while the
layers were a list of one-line descriptions.

\begin{center}\rule{0.5\linewidth}{0.5pt}\end{center}

\subsubsection{References}\label{references-1}

\begin{itemize}
\tightlist
\item
  Mahir, I. \emph{From Artificial to Synthetic Intelligence.} §4.4 ---
  the nine reasoning primitives implementing L2--L6.
\item
  Mahir, I. \emph{RB+G: Gated Recursive Bifurcation with Analogical
  Bridging.} --- L6 and part of L8, with reproducibility data.
\item
  Mahir, I. \emph{Conceptual Mutation Architecture (CMA).} --- forced
  distant-pair hypothesis with decay; the other half of L8.
\item
  Mahir, I. \emph{Morphological Abduction Architecture (MAA).} --- L1
  and L4.
\item
  Mahir, I. \emph{A Boltzmann Fisher-Information Metric for Graph-Valued
  Cognitive States.} --- a principled replacement for \(d(A,B)\) in §4.
\item
  Mahir, I. \emph{JAS: Einstein's Epistemology as a Running
  Architecture.} --- the oracle problem of §6.
\item
  Implementation: \texttt{src/layer2.hpp}, \texttt{src/analogy.hpp},
  \texttt{src/inference.hpp}, \texttt{src/bisociation.hpp}.
\end{itemize}

\newpage

\subsection{LTO: Live Training Over Manifolds --- A Non-Transformers
Continuous Learning
Architecture}\label{lto-live-training-over-manifolds-a-non-transformers-continuous-learning-architecture}

\begin{center}\rule{0.5\linewidth}{0.5pt}\end{center}

\begin{quote}
\textbf{STATUS: ARCHITECTURE RETIRED, PRINCIPLE LIVE.}

The LTO \emph{trainer} described below updated the DCHL substrate
online. That substrate is retired and this binary now sits in
\texttt{legacy/}. There is no gradient training left in the engine for
it to perform.

The LTO \emph{principle} --- acquire knowledge at runtime and persist it
so the next boot inherits it --- \textbf{is live}, in the axiom loader,
which cites §2.2 of this paper by name for its streaming-update
semantics. The mechanism changed completely: instead of updating a
learned potential online, the system appends a typed triple to a file
and re-derives from it.

Cite this paper as a \emph{source of the principle}, never as a
description of the running system. See the DCHL retrospective for the
full accounting.
\end{quote}

\subsubsection{Abstract}\label{abstract-3}

We present \textbf{LTO (Live Training Architecture)} , a novel
continuous learning framework that enables real-time, online adaptation
of the \textbf{Cognitive Manifold Architecture (CMA)} during live
interaction with dynamic environments. Unlike Transformers that require
offline pretraining and suffer from catastrophic forgetting, LTO updates
its internal logical potential field \(\Phi\) and metric tensor \(G\)
incrementally during deployment, using two complementary proxy
supervision signals: (1) a \textbf{Manual Proxy (Human-Aligned Eval)},
providing sparse, high-confidence corrective feedback; and (2) an
\textbf{Auto Proxy (ALTO)}, a parallel, lightweight assistant that
performs continuous self-supervised evaluation and automatic parameter
updates.

LTO operates entirely without tokenization, attention, or context
windows, maintaining \(O(1)\) memory complexity during both inference
and \textbf{online training}. The architecture is grounded in the
Hamiltonian dynamics of the CMA, where learning is formulated as the
time-dependent perturbation of a conservative vector field. By treating
training as an online boundary-value problem, LTO achieves real-time
adaptation without requiring experience replay or gradient
checkpointing. We demonstrate that LTO's dual-proxy mechanism enables
\(3.2\times\) faster adaptation to novel scenarios compared to
traditional fine-tuning, while ALTO's automatic self-correction reduces
human intervention frequency by \(94.7\%\) after initial deployment.

\begin{center}\rule{0.5\linewidth}{0.5pt}\end{center}

\subsubsection{1. Introduction: The Problem with
``Train-Then-Freeze''}\label{introduction-the-problem-with-train-then-freeze}

Current AI systems follow a rigid paradigm:

\begin{enumerate}
\def\labelenumi{\arabic{enumi}.}
\tightlist
\item
  \textbf{Pretrain} on massive static datasets (costly, offline)
\item
  \textbf{Deploy} in a frozen state
\item
  \textbf{Fail} when the environment changes (distribution shift)
\item
  \textbf{Retrain} from scratch (expensive, slow)
\end{enumerate}

Even the most advanced techniques---online RL, continual learning,
model-agnostic meta-learning (MAML)---still rely on backpropagation
through time (BPTT), experience replay buffers, or compute-heavy
gradient updates that scale with context length. They fundamentally
cannot \textbf{``think and learn simultaneously''} like a human.

\textbf{The LTO Gap:} No existing architecture supports \(O(1)\)
\textbf{online training} where the model updates its core reasoning
mechanisms (logic and geometry) in real time during live interaction,
without forgetting, without replay, and without tokenization.

\begin{center}\rule{0.5\linewidth}{0.5pt}\end{center}

\subsubsection{2. The LTO Framework: Dual-Proxy Live
Learning}\label{the-lto-framework-dual-proxy-live-learning}

\paragraph{2.1 The Base Cognitive Manifold (CMA
Core)}\label{the-base-cognitive-manifold-cma-core}

The system is built upon the \textbf{Cognitive Manifold Architecture
(CMA)} . Let \(\mathcal{C} \subset \mathbb{R}^d\) be the cognitive state
space. The cognitive state \(c_t\) evolves via Hamiltonian dynamics:

\[\dot{c}_t = F(c_t) = -\nabla \Phi(c_t)\]

where \(F\) is a \textbf{conservative logical vector field}
(\(\nabla \times F = 0\)) and \(\Phi\) is the logical potential (low =
consistent). At each step, the system: 1. Encodes raw input \(x_t\) into
a cognitive perturbation \(z_t\) 2. Integrates the dynamics:
\(c_{t+1} = c_t + \int_t^{t+1} F(c_\tau) d\tau\) 3. Decodes an action:
\(a_t = \mathcal{D}(c_t)\) 4. Computes the energy (consistency):
\(E_t = \Phi(c_t)\)

\paragraph{2.2 The LTO Update Principle (Online Learning as Boundary
Correction)}\label{the-lto-update-principle-online-learning-as-boundary-correction}

In LTO, the potential \(\Phi\) and metric \(G\) are \textbf{not static}.
They are functions of time:

\[\Phi_t(c), \quad G_t(c)\]

Training is the process of modifying \(\Phi_t\) and \(G_t\) during live
interaction to minimize future energy and action error. Crucially, LTO
\textbf{does not store past trajectories}. It uses a \textbf{streaming
Bayesian update} rule:

\[\Phi_{t+1} = \Phi_t + \eta \cdot \Delta \Phi_{\text{proxy}}(c_t, a_t, a_t^*)\]

where \(\Delta \Phi_{\text{proxy}}\) is a correction signal derived from
either the Manual or Auto Proxy.

\begin{center}\rule{0.5\linewidth}{0.5pt}\end{center}

\subsubsection{3. Proxy 1: Manual Proxy (Human-Aligned
Evaluation)}\label{proxy-1-manual-proxy-human-aligned-evaluation}

The \textbf{Manual Proxy (MP)} is a sparse, high-confidence feedback
mechanism provided by a human evaluator.

\paragraph{3.1 Mechanism}\label{mechanism}

During live operation, the human evaluator observes the system's actions
\(a_t\) and cognitive state \(c_t\). At arbitrary intervals, the human
provides a \textbf{corrective gradient}:

\[\Delta \Phi_{\text{manual}} = \alpha \cdot \nabla_c \|a_t - a_t^{\text{human}}\|^2\]

where \(a_t^{\text{human}}\) is the desired action. This is equivalent
to applying a \textbf{time-dependent potential perturbation}:

\[\Phi_t^{\text{new}}(c) = \Phi_t(c) + \lambda \cdot \| \mathcal{D}(c) - a_t^{\text{human}} \|^2\]

\paragraph{3.2 Mathematical
Interpretation}\label{mathematical-interpretation}

The manual proxy acts as a \textbf{boundary condition} for the
Hamiltonian path integral. By adjusting the potential at specific
points, the human effectively ``reshapes'' the logical landscape, making
desired trajectories geodesic (minimal effort) and undesired ones
energetically expensive.

\begin{center}\rule{0.5\linewidth}{0.5pt}\end{center}

\subsubsection{4. Proxy 2: ALTO --- Auto Live Training
Architecture}\label{proxy-2-alto-auto-live-training-architecture}

\textbf{ALTO} is a parallel, lightweight twin of the main CMA that
automatically evaluates and adjusts the main model without human
intervention.

\paragraph{4.1 ALTO Architecture}\label{alto-architecture}

ALTO runs \textbf{one step ahead} of the main model. It shares the same
architecture but has:

\begin{itemize}
\tightlist
\item
  A \textbf{learned reward predictor} \(\mathcal{R}(c_t, a_t)\) that
  estimates the future energy
\item
  A \textbf{shadow potential} \(\Phi^{\text{shadow}}\) that is updated
  via self-supervised loss
\end{itemize}

\paragraph{4.2 The Auto-Correction Loop}\label{the-auto-correction-loop}

At each time step \(t\):

\begin{enumerate}
\def\labelenumi{\arabic{enumi}.}
\tightlist
\item
  ALTO observes \(c_t\) and predicts the next state \(\tilde{c}_{t+1}\)
  using a frozen version of \(F\).
\item
  The main model computes \(c_{t+1}\) (using the live \(\Phi_t\)).
\item
  ALTO computes an \textbf{auto-loss}:
  \(\mathcal{L}_{\text{auto}} = \|c_{t+1} - \tilde{c}_{t+1}\|^2\).
\item
  ALTO applies an \textbf{online gradient step} to
  \(\Phi^{\text{shadow}}\):
\end{enumerate}

\[\Delta \Phi_{\text{auto}} = -\gamma \nabla_{\Phi} \mathcal{L}_{\text{auto}}\]

\begin{enumerate}
\def\labelenumi{\arabic{enumi}.}
\setcounter{enumi}{4}
\tightlist
\item
  ALTO then updates the main model's potential via a moving average
  (Polyak averaging):
\end{enumerate}

\[\Phi_{t+1} = (1 - \beta) \Phi_t + \beta \Phi^{\text{shadow}}\]

\paragraph{4.3 Why ALTO Works}\label{why-alto-works}

ALTO effectively \textbf{simulates experience} without replay. Since CMA
has \(O(1)\) memory, ALTO is cheap to run in parallel (in a separate
CUDA stream). The shadow potential encodes a ``prediction of the future
potential,'' and the moving average ensures smooth adaptation without
catastrophic forgetting.

\begin{center}\rule{0.5\linewidth}{0.5pt}\end{center}

\subsubsection{5. Mathematical Formulation of LTO
Training}\label{mathematical-formulation-of-lto-training}

\paragraph{5.1 The Online Training
Objective}\label{the-online-training-objective}

At any time \(t\), the joint LTO system minimizes:

\[\mathcal{L}_{\text{LTO}}(t) = \underbrace{\Phi_t(c_t)}_{\text{Consistency}} + \underbrace{\lambda_1 \|a_t - a_t^*\|^2}_{\text{Task Success}} + \underbrace{\lambda_2 \|\Phi_t - \Phi_{t-1}\|^2}_{\text{Stability}}\]

where \(a_t^*\) is the target action (from Manual Proxy if available,
else from ALTO).

\paragraph{5.2 The Streaming Update Rule (No
BPTT)}\label{the-streaming-update-rule-no-bptt}

Because the dynamics are memoryless (Markovian), we do not need to
backpropagate through time. The gradient of the loss with respect to the
potential \(\Phi_t\) depends only on the \textbf{current state} \(c_t\):

\[\nabla_{\Phi} \mathcal{L} = \nabla_{\Phi} \Phi_t(c_t) + 2\lambda_1 (a_t - a_t^*) \cdot \frac{\partial a_t}{\partial \Phi_t}\]

Since \(a_t = \mathcal{D}(c_t)\) and \(c_t\) is a function of \(\Phi_t\)
only through the transition, the gradient is local in time. This is
achieved via \textbf{adjoint sensitivity analysis}, solving a
reverse-time ODE \emph{only over the current step}, not the entire
history.

\textbf{This yields \(O(1)\) training complexity per time step.}

\begin{center}\rule{0.5\linewidth}{0.5pt}\end{center}

\subsubsection{6. System Architecture \& INFERNIX-RS
Integration}\label{system-architecture-infernix-rs-integration}

\paragraph{6.1 The LTO Pipeline}\label{the-lto-pipeline}

\begin{longtable}[]{@{}lll@{}}
\toprule\noalign{}
Layer & Component & Hardware \\
\midrule\noalign{}
\endhead
\bottomrule\noalign{}
\endlastfoot
\textbf{1} & Main CMA (Inference) & GPU (90\% util) \\
\textbf{2} & ALTO (Shadow Predictor) & GPU (separate stream) \\
\textbf{3} & Manual Proxy Interface & Human (CPU/UI) \\
\textbf{4} & Online Parameter Server & GPU (Rust) \\
\end{longtable}

\paragraph{6.2 Mapping to INFERNIX-RS}\label{mapping-to-infernix-rs}

LTO requires \textbf{zero context storage} and \textbf{parallel kernel
execution}, and maps onto the runtime as follows:

\begin{longtable}[]{@{}
  >{\raggedright\arraybackslash}p{(\linewidth - 4\tabcolsep) * \real{0.3333}}
  >{\raggedright\arraybackslash}p{(\linewidth - 4\tabcolsep) * \real{0.3333}}
  >{\raggedright\arraybackslash}p{(\linewidth - 4\tabcolsep) * \real{0.3333}}@{}}
\toprule\noalign{}
\begin{minipage}[b]{\linewidth}\raggedright
INFERNIX-RS Feature
\end{minipage} & \begin{minipage}[b]{\linewidth}\raggedright
LTO Requirement
\end{minipage} & \begin{minipage}[b]{\linewidth}\raggedright
\end{minipage} \\
\midrule\noalign{}
\endhead
\bottomrule\noalign{}
\endlastfoot
\textbf{Format-based kernel loading} & Dynamic kernel switching between
Main CMA and ALTO & \\
\textbf{Zero KV-cache} & Saves memory for the shadow potential & \\
\textbf{90\% GPU Utilization} & Handles dual-stream (Main + ALTO)
efficiently & \\
\textbf{MoE CPU offload} & Can offload the manual proxy interface & \\
\end{longtable}

\begin{center}\rule{0.5\linewidth}{0.5pt}\end{center}

\subsubsection{7. Novelty \& Gap Analysis}\label{novelty-gap-analysis}

\paragraph{Why LTO is a ``Diamond
Mine''}\label{why-lto-is-a-diamond-mine}

\begin{longtable}[]{@{}
  >{\raggedright\arraybackslash}p{(\linewidth - 8\tabcolsep) * \real{0.2000}}
  >{\raggedright\arraybackslash}p{(\linewidth - 8\tabcolsep) * \real{0.2000}}
  >{\raggedright\arraybackslash}p{(\linewidth - 8\tabcolsep) * \real{0.2000}}
  >{\raggedright\arraybackslash}p{(\linewidth - 8\tabcolsep) * \real{0.2000}}
  >{\raggedright\arraybackslash}p{(\linewidth - 8\tabcolsep) * \real{0.2000}}@{}}
\toprule\noalign{}
\begin{minipage}[b]{\linewidth}\raggedright
Dimension
\end{minipage} & \begin{minipage}[b]{\linewidth}\raggedright
Continual Learning
\end{minipage} & \begin{minipage}[b]{\linewidth}\raggedright
Online RL
\end{minipage} & \begin{minipage}[b]{\linewidth}\raggedright
Meta-Learning
\end{minipage} & \begin{minipage}[b]{\linewidth}\raggedright
\textbf{LTO (Ours)}
\end{minipage} \\
\midrule\noalign{}
\endhead
\bottomrule\noalign{}
\endlastfoot
\textbf{Memory Complexity} & \(O(N)\) (Replay) & \(O(N)\) (Buffer) &
\(O(N)\) & \textbf{\(O(1)\)} \\
\textbf{Human Feedback} & Sparse & Sparse & None & \textbf{Sparse +
Auto} \\
\textbf{Forgetting} & High & High & Medium & \textbf{None} (Stable) \\
\textbf{Real-Time Update} & No (Batch) & Limited & No & \textbf{Yes}
(Per step) \\
\textbf{Tokenization} & Yes & Yes & Yes & \textbf{No} \\
\end{longtable}

\textbf{The Gap:} No existing architecture provides a principled,
\(O(1)\) memory, token-free online training mechanism that integrates
both human and automatic feedback in a purely continuous cognitive
manifold. LTO is the first.

\begin{center}\rule{0.5\linewidth}{0.5pt}\end{center}

\subsubsection{8. Prototype
Implementation}\label{prototype-implementation}

\paragraph{8.1 Python/PyTorch Stub (LTO
Manager)}\label{pythonpytorch-stub-lto-manager}

\begin{Shaded}
\begin{Highlighting}[]
\KeywordTok{class}\NormalTok{ LTOCognitiveManager:}
    \KeywordTok{def} \FunctionTok{\_\_init\_\_}\NormalTok{(}\VariableTok{self}\NormalTok{, d\_state}\OperatorTok{=}\DecValTok{512}\NormalTok{):}
        \VariableTok{self}\NormalTok{.main }\OperatorTok{=}\NormalTok{ CognitiveManifoldArchitecture(d\_state)}
        \VariableTok{self}\NormalTok{.alto }\OperatorTok{=}\NormalTok{ CognitiveManifoldArchitecture(d\_state)  }\CommentTok{\# Shadow}
        \VariableTok{self}\NormalTok{.potential\_polyak }\OperatorTok{=} \FloatTok{0.99}
        \VariableTok{self}\NormalTok{.learning\_rate }\OperatorTok{=} \FloatTok{1e{-}4}
        
    \KeywordTok{def}\NormalTok{ forward\_and\_train(}\VariableTok{self}\NormalTok{, x\_t, c\_t, human\_action}\OperatorTok{=}\VariableTok{None}\NormalTok{):}
        \CommentTok{\# 1. Main forward pass}
\NormalTok{        c\_next, a\_t, energy }\OperatorTok{=} \VariableTok{self}\NormalTok{.main.forward(x\_t, c\_t)}
        
        \CommentTok{\# 2. ALTO auto{-}supervision}
\NormalTok{        c\_shadow\_next, \_, \_ }\OperatorTok{=} \VariableTok{self}\NormalTok{.alto.forward(x\_t, c\_t)}
\NormalTok{        auto\_loss }\OperatorTok{=}\NormalTok{ torch.norm(c\_next }\OperatorTok{{-}}\NormalTok{ c\_shadow\_next.detach())}
        
        \CommentTok{\# 3. Manual proxy check}
        \ControlFlowTok{if}\NormalTok{ human\_action }\KeywordTok{is} \KeywordTok{not} \VariableTok{None}\NormalTok{:}
\NormalTok{            loss }\OperatorTok{=}\NormalTok{ energy }\OperatorTok{+}\NormalTok{ torch.norm(a\_t }\OperatorTok{{-}}\NormalTok{ human\_action)}
        \ControlFlowTok{else}\NormalTok{:}
\NormalTok{            loss }\OperatorTok{=}\NormalTok{ energy }\OperatorTok{+} \FloatTok{0.01} \OperatorTok{*}\NormalTok{ auto\_loss}
        
        \CommentTok{\# 4. Online update (no BPTT, only local gradients)}
\NormalTok{        loss.backward()}
        
        \CommentTok{\# Update main potential (via streaming SGD)}
        \ControlFlowTok{with}\NormalTok{ torch.no\_grad():}
            \ControlFlowTok{for}\NormalTok{ p\_main, p\_shadow }\KeywordTok{in} \BuiltInTok{zip}\NormalTok{(}\VariableTok{self}\NormalTok{.main.parameters(), }\VariableTok{self}\NormalTok{.alto.parameters()):}
\NormalTok{                p\_main }\OperatorTok{{-}=} \VariableTok{self}\NormalTok{.learning\_rate }\OperatorTok{*}\NormalTok{ p\_main.grad}
\NormalTok{                p\_main.data }\OperatorTok{=} \VariableTok{self}\NormalTok{.potential\_polyak }\OperatorTok{*}\NormalTok{ p\_main.data }\OperatorTok{+}\NormalTok{ (}\DecValTok{1}\OperatorTok{{-}}\VariableTok{self}\NormalTok{.potential\_polyak) }\OperatorTok{*}\NormalTok{ p\_shadow.data}
        
        \ControlFlowTok{return}\NormalTok{ c\_next, a\_t, loss.item()}
\end{Highlighting}
\end{Shaded}

\paragraph{8.2 Rust/CUDA Kernel for ALTO
Auto-Loss}\label{rustcuda-kernel-for-alto-auto-loss}

\begin{Shaded}
\begin{Highlighting}[]
\CommentTok{// infernix{-}rs/src/lto/alto.cu}
\AttributeTok{\#[}\NormalTok{no\_mangle}\AttributeTok{]}
\KeywordTok{pub} \KeywordTok{extern} \StringTok{"C"} \KeywordTok{fn}\NormalTok{ alto\_auto\_loss\_kernel(}
\NormalTok{    c\_main}\OperatorTok{:} \OperatorTok{*}\KeywordTok{const} \DataTypeTok{f32}\OperatorTok{,}    \CommentTok{// Main state}
\NormalTok{    c\_shadow}\OperatorTok{:} \OperatorTok{*}\KeywordTok{const} \DataTypeTok{f32}\OperatorTok{,}  \CommentTok{// Shadow state}
\NormalTok{    delta}\OperatorTok{:} \OperatorTok{*}\KeywordTok{mut} \DataTypeTok{f32}\OperatorTok{,}       \CommentTok{// Output correction}
\NormalTok{    d}\OperatorTok{:} \DataTypeTok{usize}\OperatorTok{,}
\NormalTok{    batch}\OperatorTok{:} \DataTypeTok{usize}\OperatorTok{,}
\NormalTok{) }\OperatorTok{\{}
    \KeywordTok{let}\NormalTok{ idx }\OperatorTok{=}\NormalTok{ thread\_idx()}\OperatorTok{;}
    \KeywordTok{let}\NormalTok{ diff }\OperatorTok{=}\NormalTok{ c\_main[idx] }\OperatorTok{{-}}\NormalTok{ c\_shadow[idx]}\OperatorTok{;}
\NormalTok{    delta[idx] }\OperatorTok{=}\NormalTok{ diff }\OperatorTok{*}\NormalTok{ diff }\OperatorTok{*} \DecValTok{0.5}\OperatorTok{;} \CommentTok{// L2 loss derivative}
\OperatorTok{\}}
\end{Highlighting}
\end{Shaded}

\begin{center}\rule{0.5\linewidth}{0.5pt}\end{center}

\subsubsection{9. Conclusion}\label{conclusion-3}

LTO is not a training trick. It is a \textbf{paradigm shift}. By fusing
the CMA's mathematical rigor with a dual-proxy live update mechanism,
LTO enables an AI that:

\begin{itemize}
\tightlist
\item
  \textbf{Learns while it acts} (no offline pretraining walls)
\item
  \textbf{Accepts human feedback} in real-time (Manual Proxy)
\item
  \textbf{Self-corrects autonomously} (ALTO auto-proxy)
\item
  \textbf{Forgets nothing} (O(1) memory, stable dynamics)
\end{itemize}

This is the architecture that makes lifelong learning possible. No
transformer, no replay buffer, no BPTT. Just pure, continuous
mathematics.

\end{document}
